\chapter{Maximum Likelihood Methods}\label{ch:MLEmethods}

\section{The Model}\label{sec:Model}

An array of \(R \in \mathbb{Z}_{\geq2}\) elements receives, from a single far-field source, a stream of photons during an \textbf{observation window }of length \(W\in\mathbb{R}_{\geq0}\).
We assume that the antennas in the array use an optical detector that provides the number of photon counts per given time interval.

\iffalse
\noindent
We allow every element \(r\) within the array to have a individual \textbf{timestamp resolution} \(\Delta_r\in \mathbb{R}_{\geq0}\), but we assume a common \textbf{time slot} of length \(T_s\in \mathbb{R}_{\geq0}\) across all elements, implicitly assuming that \(T_s\) is an integer multiple of \(\Delta_r\) \(\forall r\in\{0,\ldots,R-1\}\).
\fi

\noindent
The continuous time axis \(t \in \mathbb{R}\) is therefore divided into a sequence of time slots: \((iT_s,(i+1)T_s)\) with \(i \in \mathbb{Z}_{\geq0}\). Within each slot interval \(i\), \(x_r[i]\) photons are detected by element \(r\), therefore, the number of photons arriving can be described by a discrete sequence of whole numbers \((x_r[0],x_r[1],x_r[2],\ldots)\).

\noindent
In order to evaluate the TDOA \(\tau_{jk}\) between two elements in the array, the transmitter uses a signaling scheme, and the receivers implement a strategy to estimate such a delay. 

\noindent
In~\Cref{sec:Single-Pulse-Delay-Estimation} we are going to look at a \textbf{single pulse} MLE method, in~\Cref{sec:Repeated-Pulse-Delay-Estimation} we are going to consider a \textbf{repeated pulse} scheme and in~\Cref{sec:Known-PPM-sequence-Delay-Estimation} we are going to generalize the problem to a \textbf{known PPM sequence} with inter-symbol guard times.
